\begin{mistake}{2026-04-13}{03}

\begin{problemcontent}
设 \(A, B\) 均为 \(n\) 阶矩阵，且 \(AB = A + B\)，则
(1) 若 \(A\) 可逆，则 \(B\) 可逆.
(2) 若 \(B\) 可逆，则 \(A+B\) 可逆.
(3) 若 \(B\) 可逆，则 \(A\) 可逆.
(4) \(A-E\) 恒可逆.
上述命题中，正确的命题共有（\quad）
(A) 1个. \quad (B) 2个. \quad (C) 3个. \quad (D) 4个.
\end{problemcontent}

\begin{solutioncontent}
由 \(AB = A + B\) 移项得 \(AB - A - B = 0\)。
两边同加 \(E\)，提取公因式配凑得 \((A-E)(B-E) = E\)。
根据方阵可逆定义，可知 \(A-E\) 与 \(B-E\) 均可逆，故命题 (4) 正确，且 \(|A-E| \neq 0, |B-E| \neq 0\)。

由 \(AB = A + B\)，分别提取公因式 \(A\) 和 \(B\)：
\(B = A(B-E) \implies |B| = |A| \cdot |B-E|\)
\(A = (A-E)B \implies |A| = |A-E| \cdot |B|\)
因为 \(|B-E| \neq 0\) 且 \(|A-E| \neq 0\)，可知 \(|A| \neq 0 \iff |B| \neq 0\)。即 \(A\) 可逆等价于 \(B\) 可逆，故命题 (1)、(3) 均正确。

若 \(B\) 可逆，由上述结论知 \(A\) 亦可逆，此时 \(|A| \neq 0\) 且 \(|B| \neq 0\)。
对原式两边取行列式得 \(|A+B| = |AB| = |A| \cdot |B| \neq 0\)，因此 \(A+B\) 可逆，故命题 (2) 正确。

综上所述，4 个命题均正确，选 (D)。
\end{solutioncontent}

\mistakeknowledge{
\begin{itemize}
 \item \textbf{核心公式/定理}：方阵乘积的行列式 \(|AB|=|A| \cdot |B|\)；矩阵配凑法 \((A-E)(B-E)=E\)。
 \item \textbf{易错点}：判断 \(A+B\) 可逆时，易陷入直接凑逆矩阵的误区。应利用已知条件 \(A+B=AB\) 将和转化为积，再取行列式判断是否为零。
 \item \textbf{关联知识}：方阵可逆的充要条件为 \(|A| \neq 0\)；若 \(n\) 阶方阵满足 \(MN=E\)，则 \(M, N\) 均可逆。
\end{itemize}}

\end{mistake}
